As already stated, the forkserver lies inside a \gls{SO} that is injected as dynamic dependence to the binary.
Moreover, as seen in \cref{sec:optimization}, the pointer to the forkserver function is initialized 
(togheter with the \gls{SHM} pointer) in the init function (declared as constructor and executed before the target binary).
During the execution of the ``init'' function the shared object access two environmental variables:
\begin{itemize}
    \item \textit{\_\_AFL\_SHM\_ID} : represent the \gls{SHM} ID setted by the fuzzer (as seen in \cref{sec:state_of_the_art}). 
    \item \textit{\_\_AFL\_PLUSPLUS} : this is used to understand which version of the AFL fuzzer is being used. 
        The tool supports both AFL and AFL++ communication protocol, so this environmental variable is used to choose between the two.
\end{itemize}
Depending on the value of these two environmental variable, 3 different forkserver functions could be used, 
as can be seen in 
\cref{tab:forkserver_modes}
\begin{table}[ht]
    \centering
    \caption{Different execution modes of the forkserver based on environmental variables value}
    \label{tab:forkserver_modes}
    \begin{tabular}{c|c|c}
    \hline
    \textbf{\_\_AFL\_SHM\_ID} & \textbf{\_\_AFL\_PLUSPLUS} & forkserver execution mode\\
    \hline
    NULL & NULL & \textit{standalone mode}\\
    \hline
    NULL & \textit{any value} & \textit{standalone mode}\\
    \hline
    NOT NULL & NULL & \textit{afl++ mode}\\
    \hline
    NOT NULL & 0 & \textit{afl mode}\\
    \hline
    NOT NULL & \textit{any other value} & \textit{afl++ mode}\\
        
    \end{tabular}
\end{table}
\begin{itemize}
    \item \textbf{standalone mode} : this is a modality in which the fuzzer is not present. It consent to run the binary  
    only one time (without the typical loop of the forkserver) without any communication protocol. In this case there is not 
    the \gls{SHM} so the piece of memory that is actually updated is instantiated by the forkserver with a \texttt{malloc} call.\footnote{
        this fact is completely transparent to the target binary which see the relocated memory address as it be a shared memory
    }
    \item \textbf{afl mode} : In this case there is an actual \gls{SHM} and a fuzzer that track the edge coverage.
    \item \textbf{afl++} : Same as the previous point but with a different communication protocol between the forkserver 
        and the fuzzer
\end{itemize}
